\chapter{Tema, contexto e problema}

\section{Tema selecionado}
\textbf{Tema oficial:} \placeholder{nome e número oficial do tema}

\instruction{Copie o nome oficial do tema. Não substitua o título do cenário por um nome inventado pelo grupo.}

\section{Contexto de operação}
\begin{tabularx}{\textwidth}{p{0.28\textwidth} Y}
\toprule
Campo & Preenchimento\\
\midrule
Ambiente ou processo & \field{onde o problema ocorre}\\
Público ou equipe afetada & \field{quem executa ou sofre o processo}\\
Situação atual & \field{como a atividade é realizada hoje}\\
Entradas disponíveis & \field{dados, sinais, imagens ou eventos}\\
Saídas necessárias & \field{decisão, alerta, registro ou ação}\\
Evidência do problema & \field{observação, fonte, medição ou hipótese}\\
\bottomrule
\end{tabularx}

\section{Condições de operação}
\begin{tabularx}{\textwidth}{p{0.28\textwidth} Y Y}
\toprule
Condição & Valor ou hipótese & Verificação/limite\\
\midrule
Temperatura, poeira, vibração ou ruído & \field{condição relevante} & \field{fonte ou ensaio}\\
Energia e montagem física & \field{alimentação, fixação e restrição} & \field{teste}\\
Conectividade & \field{topologia, disponibilidade e pontos cegos} & \field{perda de rede e contingência}\\
Operação offline & \field{dados armazenados localmente ou não aplicável} & \field{critério de recuperação}\\
\bottomrule
\end{tabularx}

\section{Dor ou problema específico}
\instruction{Use a fórmula: “No contexto de [ambiente], [público/processo] enfrenta [problema observável], que causa [consequência]. O projeto pretende [verbo: detectar, medir, reduzir ou registrar] [efeito], nas condições [limites].”}

\campoGrande{Descreva a situação-problema, seus atores, a falha observável e a consequência. Não atribua custo, ganho ou risco sem fonte.}

\section{Validação da necessidade}
\begin{tabularx}{\textwidth}{p{0.25\textwidth} Y Y}
\toprule
Campo & Preenchimento & Estado\\
\midrule
Fonte ou participantes & \field{fonte, processo ou atores consultados} & \field{estado}\\
Método de confirmação & \field{observação, entrevista, medição ou documento} & \field{estado}\\
Amostra ou recorte & \field{quantidade, período ou condição} & \field{estado}\\
Limitação da evidência & \field{o que a fonte não permite concluir} & \field{estado}\\
Critério de aceitação operacional & \field{quem considera o resultado útil} & \field{estado}\\
\bottomrule
\end{tabularx}

\section{Objetivo e resultado pretendido}
\instruction{Descreva o artefato ou comportamento que a solução tentará produzir. Separe resultado pretendido, resultado medido e resultado não alcançado.}

\textbf{Objetivo geral:} \field{verbo + objeto + condição de uso}

\textbf{Resultado pretendido:} \field{o que deverá existir ou acontecer ao final}

\textbf{Resultado medido nesta entrega:} \field{não medido / preencher somente se houver ensaio}

\section{Proposta de valor}
Explique como o resultado pretendido responde à dor. Impacto financeiro, desempenho operacional, economia ou segurança deve ser citado como expectativa, fonte ou medição, nunca como consequência garantida.

\section{Roteiro de comunicação}
\begin{tabularx}{\textwidth}{p{0.25\textwidth} Y}
\toprule
Parte & Preenchimento\\
\midrule
Problema & \field{qual dor deve prender a atenção da banca}\\
Solução & \field{como a arquitetura responde ao problema}\\
Evidências & \field{quais medições, PoCs ou artefatos sustentam a proposta}\\
Continuidade & \field{qual extensão depende de validação futura}\\
\bottomrule
\end{tabularx}

\instruction{Traduza o jargão em função e benefício. Não apresente impacto esperado como resultado comprovado.}
